\section{Introduction}
\label{sec:intro}

Computational design and engineering problems rarely end with a single reusable design. In
emergency-facility planning, land availability, district constraints or budget revisions can make
a nominally optimal layout unusable after the analysis is complete. A more useful decision-support
optimiser returns several layouts that are nearly indistinguishable in objective value but clearly
different in geography. The same need for distinct solutions of comparable quality appears in
engineering design, scheduling and path planning. A single Pareto set cannot provide that
flexibility.

Multimodal multi-objective optimisation (MMOP) formalises this need: the goal is to recover
not one Pareto set but the \emph{multiple, geometrically distinct decision-space solution
sets} that all map to the (near-)same Pareto front~\parencite{yue2018moring,yue2019mmf}. Good MMOP
solutions therefore support revision without restarting the search. Poor solutions can present a
frontier while still hiding a single brittle decision pattern. The central difficulty is that the
selection pressure needed for convergence also removes the decision-space diversity that MMOP is
meant to preserve.

\paragraph{The fundamental obstacle} This is the \emph{convergence/diversity trade-off in
decision space}. An optimiser that drives hard toward the front can collapse the population onto a
single decision region. An optimiser that spreads solutions aggressively can lose objective
quality. State-of-the-art MMOP methods usually manage this trade-off rather than remove it.
CPDEA~\parencite{liu2020cpdea} adds a convergence-penalised decision-space density, and
MMEA-WI~\parencite{li2021mmeawi} folds decision diversity into a weighted indicator; both
\emph{fuse} convergence and diversity into a single scalar fitness, so buying more
decision-space coverage necessarily spends some convergence. Classic niching methods
(MO\_Ring\_PSO\_SCD~\parencite{yue2018moring}, DN-NSGA-II~\parencite{liang2016dnnsga2},
Omni-optimizer~\parencite{deb2008omni}) sit at different points of the same curve. The open
question is therefore not how to weight the two objectives better, but whether decision-space
diversity can be added at only a \emph{small, bounded} convergence cost instead of the full
trade-off these methods pay.

\paragraph{Our key insight: \emph{where} the diversity term enters the pipeline} We answer
this question by moving the decision-space signal to a later point in selection. A generational
MMOEA first sorts solutions by dominance and then chooses among the members of the splitting
front. Prior state-of-the-art methods let decision-space diversity influence the first step, either through
convergence-penalised density or through a weighted indicator. EARS-MMOEA
(\emph{Equivalence-Aware, Rarity-boosted, Structure-preserving} MMOEA) leaves the dominance sort
untouched. It applies diversity only inside the splitting front:
$D(i)=E(i)\,(1+\beta S(i))$, where $E$ is an equivalence-aware special crowding distance and
$S$ is a decision-space sparsity bonus.

This placement gives the paper its central mechanism. Since the bonus is non-negative and only
compares mutually non-dominated candidates, it preserves front precedence. A dominated solution
cannot overtake a solution from a better front. The claim is not a global convergence theorem;
it is a selection-layer separation that we measure directly. Under matched $E$ and $S$, a
controlled experiment shows that within-front placement avoids most of the convergence cost of
in-sort fusion. The precise algebraic form is secondary: additive and multiplicative
within-front versions behave similarly, while moving the same sparsity signal into the sort
significantly hurts convergence.

The paper has one primary contribution and two supporting claims:
\begin{itemize}
  \item \textbf{(C1) Primary mechanism.} A pure decision-space sparsity signal should be kept out
  of the dominance sort and used as a within-front selection term. This placement adds
  decision-space coverage at a small measured convergence cost, whereas in-sort fusion spends much
  more convergence for the same signal (Section~\ref{sec:novelty}).
  \item \textbf{(C2) Algorithmic realisation.} EARS-MMOEA implements the placement in an
  equivalence-aware framework. EARS-Core isolates the selection idea, while EARS-Full adds
  archives, adaptive niching and an operator portfolio for the competitive comparisons
  (Sections~\ref{sec:method}--\ref{sec:benchmark}).
  \item \textbf{(C3) Evidence and scope.} Benchmarks, scalable families and a five-city real-map
  OpenStreetMap facility-placement case study show the value of the placement in the target regime of
  equivalent global Pareto sets. Deceptive local Pareto sets form the reported counter-regime
  (Section~\ref{sec:experiments}).
\end{itemize}

We treat threats to validity as part of the claim. Baselines use their authors' recommended
operators and are checked on MMF. The decision-space primary metrics are fixed a priori. Negative
cases are reported, including two cities where a baseline edges EARS on the core-suite rank.
The method targets multimodality from \emph{equivalent global} Pareto sets. On deceptive local
Pareto sets and weakly multimodal tasks with an essentially unique optimum, it has little or no
advantage (Section~\ref{sec:experiments}).
